%! Adding \& Subtracting Rational Expressions — Unit 4, Lesson 4.3 — Student Packet Cover Sheet
\documentclass[10pt]{article}
\usepackage{algebra2-article}
\usepackage{algebra2-boxes}
\usepackage{ltablex}
\keepXColumns

\begin{document}

% Full-bleed forest banner
\begin{tikzpicture}[remember picture, overlay]
  \fill[forest] (current page.north west)
    rectangle ([yshift=-0.9in]current page.north east);
\end{tikzpicture}
\vspace{-0.8in}

\begin{center}
  {\color{white}\textbf{\LARGE Algebra 2: Shepherd}}\\[4pt]
  {\color{forestmid}\large\textbf{Unit 4 \quad Rational Functions}}\\[6pt]
  {\color{white}\textbf{\large Lesson 4.3 \quad Adding \& Subtracting Rational Expressions}}
\end{center}

\vspace{0.22in}
\namedateperiod
\vspace{0.18in}

\begin{learningtargetbox}
\small
\begin{enumerate}[leftmargin=1.5em, itemsep=5pt]
  \item \ldots{} add and subtract rational expressions with \textbf{like denominators}, and I put the second numerator in \textbf{parentheses} every time I subtract, because the minus sign owns the whole numerator.
  \item \ldots{} build the \textbf{least common denominator} out of the \emph{factored} denominators --- each distinct factor to its highest power --- rewrite each fraction with a \textbf{building factor}, combine, and simplify.
  \item \ldots{} recognize a \textbf{complex fraction} and simplify it as a quotient of two simple fractions, because the main fraction bar is a division sign.
  \item \ldots{} state \textbf{every} restriction, reading it off the original denominators --- including the one hiding under the main fraction bar --- and say which excluded value is a \textbf{hole} and which is a \textbf{wall}.
\end{enumerate}
\end{learningtargetbox}

\vspace{0.14in}

\begin{tocbox}
\small
\renewcommand{\arraystretch}{1.5}
\begin{tabularx}{\linewidth}{@{} c l X r @{}}
\rowcolor{forestbg}
\textbf{\#} & \textbf{Component} & \textbf{Description} & \textbf{Score} \\
\midrule
1 & Warm-Up        & Spiral review: the LCD on number fractions, subtracting past a parenthesis, factoring, and simplify-and-restrict & \blank{1.2cm} \\
2 & Guided Notes   & Like denominators, the minus-sign rule, building the LCD, and complex fractions & \blank{1.2cm} \\
3 & Group Activity & \emph{Build the Denominator, Guard the Sign} --- an LCD table, an error analysis, and a complex fraction two ways, by tier & \blank{1.2cm} \\
4 & Exit Ticket    & Quick check: an LCD subtraction with a hole and a wall, plus an SOL-style item & \blank{1.2cm} \\
5 & Homework       & Like and unlike denominators, complex fractions, a round-trip model, and a design-your-own extension & \blank{1.2cm} \\
\midrule
  & \multicolumn{2}{r}{\textbf{Total}} & \blank{1.2cm} \\
\end{tabularx}
\end{tocbox}

\vspace{0.10in}

\begin{remindbox}
\small To \textbf{add or subtract} rational expressions the denominators must match, so the work is all
in building one. Factor every denominator, then build the \textbf{least common denominator} from those
factors --- each \emph{distinct} factor, raised to the \emph{highest} power it reaches in any one
denominator. Never multiply the denominators together and never add them: multiplying is only
expensive, but adding is simply wrong ($\frac1x+\frac13\neq\frac{2}{x+3}$). Multiply each fraction by
the \textbf{building factor} $\frac kk$ that supplies what it is missing, combine the numerators over
the LCD, and only then factor the numerator to see whether anything divides out. \textbf{Every time you
subtract, put the second numerator in parentheses} and distribute the minus sign --- it changes the
sign of every term inside, and forgetting it is the single most common lost point on this topic. A
\textbf{complex fraction} --- one with fractions inside it --- is a division in disguise: the main
fraction bar is a $\div$ sign, so combine the top into one fraction, combine the bottom into one
fraction, and keep--change--flip. Its restrictions come from two levels: every inner denominator, and
the whole bottom, which may never equal zero. And the shortcut worth remembering: once your denominators
are factored, \textbf{the LCD is the restriction list} --- set each of its factors to zero and you have
every excluded value, whether or not it survives into your final answer.
\end{remindbox}

\end{document}
